\documentclass[11pt,a4paper]{article}

\usepackage{amsmath,amssymb,amsthm}
\usepackage{hyperref}
\usepackage[margin=1in]{geometry}
\usepackage{xcolor}
\usepackage{booktabs}
\usepackage{array}
\usepackage{float}
\usepackage{mathtools}

\newtheorem{theorem}{Theorem}[section]
\newtheorem{conjecture}[theorem]{Conjecture}
\newtheorem{corollary}[theorem]{Corollary}
\newtheorem{lemma}[theorem]{Lemma}
\newtheorem{proposition}[theorem]{Proposition}
\newtheorem{definition}[theorem]{Definition}
\theoremstyle{remark}
\newtheorem{remark}[theorem]{Remark}

\newcommand{\PT}{\mathbb{PT}}
\newcommand{\CP}{\mathbb{CP}}
\newcommand{\C}{\mathbb{C}}
\newcommand{\Z}{\mathbb{Z}}
\newcommand{\R}{\mathbb{R}}
\newcommand{\F}{\mathbb{F}}
\newcommand{\cO}{\mathcal{O}}
\newcommand{\cC}{\mathcal{C}}
\newcommand{\cN}{\mathcal{N}}
\newcommand{\cU}{\mathcal{U}}
\newcommand{\cW}{\mathcal{W}}
\newcommand{\cG}{\mathcal{G}}
\newcommand{\cF}{\mathcal{F}}
\newcommand{\PP}{\mathcal{P}}
\newcommand{\barPP}{\overline{\mathcal{P}}}


\title{The T-Vector Theorem: Symplectic Characteristic Elements\\
       and the Kochen--Specker Obstruction\\
       {\large (Paper XII)}}


\author{Zhou-Li Chen \\ co-nlang Research}
\date{June 2026}

\begin{document}

\maketitle

\begin{abstract}
  Paper~XI~\cite{PaperXI} established the $\beta$-formula for context
  signs of Mermin pentagrams and proved algebraically that
  $\omega(r_{ij}, r_{kl}) = 1$ for every disjoint ray pair, yielding
  $\omega_{\text{total}} \equiv 1 \pmod{2}$. The parity theorem
  $\beta_{\text{sum}} \equiv 2 \pmod{4}$ was verified computationally
  for all $12{,}096$ pentagrams.

  In this paper we prove the \emph{T-vector theorem}: for any
  equiangular Mermin pentagram with 10 rays $\{r_{ij}\}$, the sum
  $T = \sum r_{ij}$ satisfies $\omega(T, r) = 1$ for every ray $r$
  (Theorem~\ref{thm:T-vector}). The proof is purely algebraic and
  $G_2(2)$-equivariant.

  We interpret $T$ as a \emph{symplectic characteristic element} of the
  10-ray system, drawing an analogy with Wu classes in algebraic
  topology (Remark~\ref{rem:wu}).

  We then prove a negative result: the Arf invariant framework for the
  KS obstruction is incompatible with $G_2(2)$ symmetry. Specifically,
  $G_2(2)$ acts transitively on $\F_2^6 \setminus \{0\}$
  (Proposition~\ref{prop:transitive}) and no $G_2(2)$-invariant
  quadratic form exists
  (Proposition~\ref{prop:no-invariant-form}). The obstruction class
  $[f_3] \in H^1(K_5, \cF)$ therefore requires a cohomological
  interpretation beyond quadratic refinements.

  We close with the mod~4 lifting problem: an algebraic derivation of
  $\beta_{\text{sum}} \equiv 2 \pmod{4}$ from the T-vector theorem
  remains open (Conjecture~\ref{conj:mod4}).
\end{abstract}

%------------------------------------------------------------------
\section{Introduction}
\label{sec:intro}
%------------------------------------------------------------------

The Kochen--Specker (KS) theorem~\cite{KochenSpecker1967} forbids
non-contextual hidden variable models of quantum mechanics. In the
three-qubit setting, Mermin's pentagram~\cite{Mermin1990} provides a
state-independent proof using 10 observables arranged in 5 contexts of
4 mutually commuting operators.

Paper~X~\cite{PaperX} established the \emph{equiangular
characterization}: every Mermin pentagram has Lagrangian configuration
with $\dim(L_i \cap L_j) = 1$ for all 10 pairs, and conversely every
equiangular $K_5$ in $W(5, \F_2)$ is a Mermin pentagram. The 10 shared
operators form a 10-point cap in $\text{PG}(5, \F_2)$ spanning the full
space $\F_2^6$.

Paper~XI~\cite{PaperXI} provided an explicit formula for context signs
via the $\beta$-formula and proved algebraically that $\omega(r_{ij},
r_{kl}) = 1$ for every disjoint ray pair. The parity theorem
$\beta_{\text{sum}} \equiv 2 \pmod{4}$ was verified computationally.

\medskip

\textbf{Main results of this paper.}
\begin{enumerate}
  \item \textbf{T-vector theorem} (Theorem~\ref{thm:T-vector}):
    For any equiangular pentagram, $T = \sum r_{ij}$ satisfies
    $\omega(T, r) = 1$ for all rays. Proof is algebraic and
    $G_2(2)$-equivariant.
  \item \textbf{Symplectic characteristic element}
    (Remark~\ref{rem:wu}): $T$ is the unique element encoding the
    symplectic structure of the ray system, analogous to the Wu class
    in algebraic topology.
  \item \textbf{Arf framework incompatibility}
    (Theorem~\ref{thm:arf-fails}): The identification $[f_3] =
    \Phi^*(\text{Arf}(q'))$ is incompatible with $G_2(2)$ symmetry.
    $G_2(2)$ acts transitively on $\F_2^6 \setminus \{0\}$ and no
    $G_2(2)$-invariant quadratic form exists.
  \item \textbf{Mod~4 lifting problem}
    (Conjecture~\ref{conj:mod4}): An algebraic derivation of
    $\beta_{\text{sum}} \equiv 2 \pmod{4}$ from the T-vector theorem
    remains open.
\end{enumerate}

\medskip

\textbf{Structure.}
Section~\ref{sec:T-vector} proves the T-vector theorem and its
corollaries. Section~\ref{sec:symplectic-char} develops the geometric
interpretation as symplectic characteristic element.
Section~\ref{sec:arf-fails} proves the incompatibility of the Arf
framework. Section~\ref{sec:mod4} formulates the mod~4 lifting problem.
Section~\ref{sec:open} lists open problems.

%------------------------------------------------------------------
\section{The T-Vector Theorem}
\label{sec:T-vector}
%------------------------------------------------------------------

\begin{definition}[T-vector]
\label{def:T}
Let $P$ be a Mermin pentagram with 10 rays
$\{r_{ij}\}_{1 \leq i < j \leq 5} \subset \F_2^6$, where $r_{ij}$ is
the unique shared operator between contexts $C_i$ and $C_j$. The
\emph{T-vector} of $P$ is
\[
  T = \bigoplus_{1 \leq i < j \leq 5} r_{ij} \in \F_2^6.
\]
\end{definition}

\begin{theorem}[T-vector theorem]
\label{thm:T-vector}
For every equiangular Mermin pentagram and every ray $r_{ab}$:
\[
  \omega(T, r_{ab}) = 1 \pmod{2}.
\]
\end{theorem}

\begin{proof}
Fix a ray $r_{ab}$. The 10 rays decompose relative to the edge $(a,b)$
of $K_5$:

\medskip

\textbf{Sharing rays} (6): $r_{ac}, r_{ad}, r_{ae}, r_{bc}, r_{bd},
r_{be}$, where $c, d, e$ are the remaining vertices. Each sharing ray
lies in the same Lagrangian as $r_{ab}$ (both belong to context $C_a$
or $C_b$), so $\omega(r_{ac}, r_{ab}) = 0$, etc.

\medskip

\textbf{Disjoint rays} (3): $r_{cd}, r_{ce}, r_{de}$ where
$\{c,d,e\} = \{1,2,3,4,5\} \setminus \{a,b\}$. By the disjoint pair
theorem (Paper~XI, Theorem~3.5), $\omega(r_{kl}, r_{ab}) = 1$ for
each disjoint pair.

\medskip

Since $T = \bigoplus_k r_k$ and $\omega$ is bilinear over $\F_2$:
\[
  \omega(T, r_{ab})
  = \sum_{k} \omega(r_k, r_{ab})
  = 6 \times 0 + 3 \times 1
  = 3 \equiv 1 \pmod{2}. \quad \square
\]
\end{proof}

\begin{corollary}
\label{cor:T-not-in-L}
$T \notin L_i$ for any Lagrangian $L_i$ of the pentagram.
\end{corollary}

\begin{proof}
If $T \in L_i$, then $\omega(T, v) = 0$ for all $v \in L_i$, including
the rays of context $C_i$. But $\omega(T, r) = 1$ for all rays by
Theorem~\ref{thm:T-vector}.
\end{proof}

\begin{corollary}
\label{cor:T-nonzero}
$T \neq 0$ for every Mermin pentagram.
\end{corollary}

\begin{remark}[Computational properties of $T$]
\label{rem:T-stats}
The following are verified computationally for all $12{,}096$ Mermin
pentagrams~\cite{PaperXIIsuppl}:
\begin{itemize}
  \item $T$ takes all 63 non-zero values in $\F_2^6$, each appearing
    exactly 192 times.
  \item The integer symplectic form $\omega_{\text{int}}(T, r) \in
    \{-3, -1, +1, +3\}$, always odd.
  \item $q(T) = 0$ for $6{,}720$ pentagrams (55.6\%) and $q(T) = 1$
    for $5{,}376$ pentagrams (44.4\%), where $q(v) = x_1 z_1 + x_2 z_2
    + x_3 z_3$ is the standard quadratic form.
\end{itemize}
\end{remark}

%------------------------------------------------------------------
\section{Symplectic Characteristic Elements}
\label{sec:symplectic-char}
%------------------------------------------------------------------

\begin{definition}[Symplectic characteristic element]
\label{def:symplectic-char}
Let $(V, \omega)$ be a symplectic vector space over $\F_2$ and
$S = \{r_1, \ldots, r_n\} \subset V$ a finite set. A
\emph{symplectic characteristic element} of $S$ is a vector $T \in V$
satisfying
\[
  \omega(T, r_i) = 1 \quad \text{for all } r_i \in S.
\]
\end{definition}

Theorem~\ref{thm:T-vector} states that for any equiangular pentagram,
the sum $T = \sum r_{ij}$ is a symplectic characteristic element of the
10-ray system.

\begin{proposition}[Uniqueness]
\label{prop:uniqueness}
For equiangular pentagrams, the symplectic characteristic element is
unique.
\end{proposition}

\begin{proof}
If $T'$ is another symplectic characteristic element, then
$\omega(T - T', r) = 0$ for all rays $r$, so $T - T' \in S^\perp$
where $S$ is the set of 10 rays. By the 10-ray cap theorem
of~\cite{PaperX}, the rays span $\F_2^6$, so $S^\perp = \{0\}$ and
$T = T'$.
\end{proof}

\begin{remark}[Analogy with Wu classes]
\label{rem:wu}
In algebraic topology, the Wu class $v \in H^*(M; \F_2)$ of a closed
manifold $M$ is characterized by
\[
  \text{Sq}^i(x) = v \cup x \quad \text{for all } x \in H^*(M; \F_2).
\]
The T-vector is the finite-field symplectic geometry analog:

\medskip

\begin{center}
\begin{tabular}{l l l}
\toprule
& \textbf{Wu class} & \textbf{T-vector} \\
\midrule
Domain & Manifold cohomology & Symplectic vector space \\
Characterizes & Steenrod operations & Symplectic pairing \\
Via & Cup product & $\omega(\cdot, \cdot)$ \\
Equation & $\text{Sq}^i(x) = v \cup x$ & $\omega(T, r) = 1$ \\
\bottomrule
\end{tabular}
\end{center}

\medskip

Both are ``universal'' elements encoding global structure through local
operations. The existence of a symplectic characteristic element
implies that the ray system has \emph{maximal non-isotropy}: the rays
cannot all lie in a Lagrangian subspace, and $T$ serves as a
``witness'' to the symplectic structure of the configuration.
\end{remark}

\begin{remark}[KS obstruction as symplectic phenomenon]
\label{rem:symplectic-ks}
The T-vector theorem reveals that the KS obstruction is fundamentally a
\emph{symplectic} phenomenon. The obstruction arises from the symplectic
pairing structure of the rays, not from the Arf invariant of a
quadratic form. This is made precise by the incompatibility result of
Section~\ref{sec:arf-fails}.
\end{remark}

%------------------------------------------------------------------
\section{Incompatibility of the Arf Framework}
\label{sec:arf-fails}
%------------------------------------------------------------------

The symplectic space $(\F_2^6, \omega)$ admits $2^6 = 64$ quadratic
refinements $q_u = q + \ell_u$ where $q(v) = x_1 z_1 + x_2 z_2 + x_3
z_3$ is the standard form and $\ell_u(v) = \omega(u, v)$. The Arf
invariant classifies these into two families:
\begin{itemize}
  \item \textbf{Even} ($\text{Arf} = 0$): 36 forms
    ($u = 0$ or $q(u) = 0$, giving $1 + 35 = 36$ choices)
  \item \textbf{Odd} ($\text{Arf} = 1$): 28 forms
    ($q(u) = 1$, giving 28 choices)
\end{itemize}

The shift formula gives $\text{Arf}(q + \ell_u) = \text{Arf}(q) +
q(u)$, so $\text{Arf}(q_u) = q(u)$ since $\text{Arf}(q) = 0$.

\begin{remark}[Correction to Paper~XI]
\label{rem:paper-xi-correction}
Paper~XI~\cite{PaperXI}, \S6.3 incorrectly stated ``the Arf invariant
is 1 for the standard quadratic form on $\F_2^6$.'' The correct value
is $\text{Arf}(q) = 0$, as computed above via $\text{Arf}(q) = \sum_k
q(e_k)q(f_k) = 0$ in any symplectic basis. Paper~XI \S6.3 has been
corrected in the Zenodo version; the present paper uses the correct
value throughout.
\end{remark}

\medskip

The original conjecture~\cite{PaperIX} was that the KS obstruction
class $[f_3] \in H^1(K_5, \cF)$ could be identified with the pullback
$\Phi^*(\text{Arf}(q'))$ for some geometrically natural quadratic form
$q'$. We now prove this is impossible.

\begin{proposition}[$G_2(2)$ transitivity, computational]
\label{prop:transitive}
$G_2(2)$ (order $12{,}096$) acts transitively on $\F_2^6 \setminus
\{0\}$. There is a single orbit of size 63.
\end{proposition}

\begin{proof}
Verified by exhaustive computation~\cite{PaperXIIsuppl}: starting from
any non-zero vector, the orbit under the two generators of $G_2(2)$
contains all 63 non-zero vectors.
\end{proof}

\begin{corollary}
\label{cor:mixes-q}
$G_2(2)$ does not preserve the standard quadratic form $q$. The orbit
of any vector contains 35 vectors with $q = 0$ and 28 with $q = 1$.
\end{corollary}

\begin{proof}
Since $G_2(2)$ is transitive on all 63 non-zero vectors, the orbit of
any vector is the entire set $\F_2^6 \setminus \{0\}$, which contains
35 vectors with $q = 0$ and 28 with $q = 1$.
\end{proof}

\begin{proposition}[No invariant quadratic form, computational]
\label{prop:no-invariant-form}
Among all 64 quadratic forms refining $\omega$, none is
$G_2(2)$-invariant.
\end{proposition}

\begin{proof}
Verified by exhaustive search~\cite{PaperXIIsuppl}: for each of the 64
forms $q_u$, there exists $g \in G_2(2)$ and $v \in \F_2^6$ such that
$q_u(g \cdot v) \neq q_u(v)$.
\end{proof}

\begin{remark}[Conceptual reason]
\label{rem:conceptual}
Proposition~\ref{prop:no-invariant-form} follows immediately from
Proposition~\ref{prop:transitive} without computation. Since $G_2(2)$
acts transitively on $\F_2^6 \setminus \{0\}$, any $G_2(2)$-invariant
function $q: \F_2^6 \to \F_2$ must be constant on non-zero vectors.
But $q(0) = 0$ for any quadratic form, so $q$ would be either
identically zero (not a refinement of $\omega$) or identically one on
non-zero vectors (which fails $q(a+b) + q(a) + q(b) = \omega(a,b)$).
\end{remark}

\begin{theorem}[No canonical Arf identification]
\label{thm:arf-fails}
There is no $G_2(2)$-canonical identification of $[f_3]$ with the Arf
invariant of a quadratic form: no $G_2(2)$-invariant quadratic form
$q'$ refining $\omega$ satisfies $\emph{Arf}(q') = 1$.
\end{theorem}

\begin{proof}
By Proposition~\ref{prop:no-invariant-form}, no $G_2(2)$-invariant
quadratic form exists at all, so no $G_2(2)$-invariant $q'$ with
$\text{Arf}(q') = 1$ can be found.

Any geometrically natural (i.e., $G_2(2)$-canonical) identification of
$[f_3]$ with an Arf invariant would require a $G_2(2)$-invariant choice
of $q'$. Since no such choice exists, the identification cannot be made
canonically.
\end{proof}

\begin{remark}[Non-canonical identifications exist]
\label{rem:noncanonical}
For any of the 28 quadratic forms with $\text{Arf} = 1$, one may
formally write $[f_3] = 1 = \text{Arf}(q')$. But this identification
is non-canonical: there is no $G_2(2)$-equivariant rule to select $q'$
among the 28 odd forms, and no single $q'$ distinguishes the KS
obstruction geometrically.
\end{remark}

\begin{remark}[No universal dual vector]
\label{rem:no-universal-u}
For each fixed $u$ with $q(u) = 1$, the vector $u$ belongs to exactly
15 Lagrangians and $u \in \bigcup_i L_i$ for exactly $4{,}800$ of the
$12{,}096$ pentagrams (39.7\%). No single $u$ covers all pentagrams.
This is consistent with Theorem~\ref{thm:arf-fails}: the ``correct''
dual vector would need to be pentagram-dependent, but no
$G_2(2)$-equivariant selection rule exists.
\end{remark}

%------------------------------------------------------------------
\section{The Mod~4 Lifting Problem}
\label{sec:mod4}
%------------------------------------------------------------------

Paper~XI established two results at different levels:
\begin{enumerate}
  \item \textbf{Mod~2} (algebraic): $\omega_{\text{total}} \equiv 1
    \pmod{2}$, proved via the disjoint pair theorem.
  \item \textbf{Mod~4} (computational): $\beta_{\text{sum}} \equiv 2
    \pmod{4}$, verified for all $12{,}096$ pentagrams.
\end{enumerate}

The T-vector theorem (Theorem~\ref{thm:T-vector}) is a consequence of
the mod~2 result. The gap between mod~2 and mod~4 remains the central
open problem.

\begin{conjecture}[Mod~4 lifting]
\label{conj:mod4}
The congruence $\beta_{\text{sum}} \equiv 2 \pmod{4}$ can be derived
algebraically from the T-vector theorem $\omega(T, r) = 1$ and the
equiangular geometry.
\end{conjecture}

\begin{remark}[Connection to the $\beta$-formula]
\label{rem:beta-connection}
The $\beta$-formula of Paper~XI gives the context sign as
$s(C_i) = (-1)^{\beta(C_i)/2}$ where
\[
  \beta(C_i) = \sum_{j<k} \omega_{\text{int}}(v_j, v_k)
\]
and $\omega_{\text{int}}$ is the \emph{integer} symplectic form. The
total parity is
\[
  \prod_{i=1}^5 s(C_i) = (-1)^{\sum_i \beta(C_i)/2}.
\]
The sum $\sum_i \beta(C_i)$ involves $\omega_{\text{int}}$ evaluated on
the 30 sharing pairs (pairs of rays sharing a $K_5$ vertex). Each
sharing pair lies in a Lagrangian, so $\omega_{\text{int}}$ is even and
$\omega_{\text{int}}/2$ is an integer.

\medskip

The quadratic form identity
\[
  q(T) + \sum_{a=1}^{10} q(r_a) \equiv \omega_{\text{total}} \equiv 1
  \pmod{2}
\]
connects the T-vector to the integer structure. Lifting this to mod~4
requires understanding the relationship between $\omega_{\text{int}}$
(integer-valued) and $\omega$ (mod~2 reduction).
\end{remark}

\begin{remark}[Parity distribution]
\label{rem:parity-dist}
Among the $12{,}096$ Mermin pentagrams, the number of negative contexts
is always odd:
\begin{center}
\begin{tabular}{l r r}
\toprule
\textbf{Parity} & \textbf{Count} & \textbf{Fraction} \\
\midrule
1-minus & $7{,}884$ & 65.2\% \\
3-minus & $4{,}104$ & 33.9\% \\
5-minus & $108$ & 0.9\% \\
\bottomrule
\end{tabular}
\end{center}
All counts are $G_2(2)$-invariant. The 5-minus pentagrams (108 total)
are the most symmetric configurations.
\end{remark}

%------------------------------------------------------------------
\section{Open Problems}
\label{sec:open}
%------------------------------------------------------------------

\begin{enumerate}
  \item \textbf{Algebraic proof of mod~4 lifting.}
    Derive $\beta_{\text{sum}} \equiv 2 \pmod{4}$ from the T-vector
    theorem and the equiangular geometry, closing the gap between the
    algebraic mod~2 result and the computational mod~4 verification.

  \item \textbf{Weil representation.}
    The $\beta$-formula may arise naturally from the Weil representation
    of $\text{Sp}(6, \F_2)$. The integer symplectic form
    $\omega_{\text{int}}$ and the phase $i^{\beta}$ suggest a connection
    to the metaplectic representation~\cite{LionVergne}.

  \item \textbf{Higher cohomological operations.}
    Since the Arf invariant framework fails
    (Theorem~\ref{thm:arf-fails}), the obstruction $[f_3]$ may
    correspond to a higher-degree cohomological operation: Steenrod
    squares, Massey products, or characteristic classes of a gerbe on
    the context complex.

  \item \textbf{$n$-qubit generalization.}
    Does the T-vector theorem generalize to $n$-qubit Mermin stars?
    For an $n$-qubit system with $N$ rays in $\F_2^{2n}$, the sum $T =
    \sum r_a$ may satisfy $\omega(T, r_a) = c$ for some constant $c$
    determined by the combinatorial structure.

  \item \textbf{Class~II geometric characterization.}
    The $G_2(2)$ orbit decomposition yields two geometric classes
    (Paper~X, Corollary~8.2). Class~II pentagrams have enhanced
    transverse obstruction ($+5.5$ pp in Type~F4). A geometric
    invariant distinguishing the classes remains to be identified.

  \item \textbf{Split Cayley Hexagon.}
    The T-vector $T$ takes all 63 non-zero values uniformly. The
    connection between $T$ and the Split Cayley Hexagon
    substructures~\cite{LevayPlanatSaniga} of $W(5, \F_2)$ is open.
\end{enumerate}

%------------------------------------------------------------------
\section*{Appendix: Rigor Self-Assessment}
\label{app:rigor}
%------------------------------------------------------------------

\begin{table}[H]
\centering
\begin{tabular}{l l l}
\toprule
\textbf{Result} & \textbf{Label} & \textbf{Verification} \\
\midrule
T-vector theorem & thm:T-vector & Algebraic \\
$T \notin L_i$ & cor:T-not-in-L & Algebraic \\
$T \neq 0$ & cor:T-nonzero & Algebraic \\
Uniqueness of $T$ & prop:uniqueness & Algebraic \\
$G_2(2)$ transitivity & prop:transitive & Computational \\
No invariant form & prop:no-invariant-form & Computational \\
Arf incompatibility & thm:arf-fails & Algebraic (from comp.\ lemmas) \\
No universal $u$ & rem:no-universal-u & Computational \\
Mod~4 lifting & conj:mod4 & Open \\
$T$-vector statistics & rem:T-stats & Computational \\
Parity distribution & rem:parity-dist & Computational \\
\bottomrule
\end{tabular}
\caption{Rigor classification of main results.}
\label{tab:rigor}
\end{table}

%------------------------------------------------------------------
\section*{Acknowledgments}
%------------------------------------------------------------------

This work continues the computational-algebraic investigation of the
symplectic geometry of Mermin pentagrams. The computations were
performed using Python with NumPy. Scripts are archived
in~\cite{PaperXIIsuppl}.

%------------------------------------------------------------------
\begin{thebibliography}{9}
%------------------------------------------------------------------

\bibitem[PaperIX]{PaperIX}
Z.-L.~Chen,
``The 3-Qubit Obstruction Ladder: Hyperbolic Embedding,
Mermin Pentagram, and Two Klein Quartic Bridges,''
\textit{Zenodo, DOI: 10.5281/zenodo.20465623} (2026).

\bibitem[KochenSpecker1967]{KochenSpecker1967}
S.~Kochen and E.~P.~Specker,
``The problem of hidden variables in quantum mechanics,''
\textit{J. Math. Mech.} \textbf{17}, 59--87 (1967).

\bibitem[LevayPlanatSaniga]{LevayPlanatSaniga}
P.~L\'evay, M.~Planat, and P.~Saniga,
``Three-qubit Mermin pentagram, the split Cayley hexagon, and
the black-hole analogy,''
\textit{arXiv:1303.5653} (2013).

\bibitem[LionVergne]{LionVergne}
G.~Lion and M.~Vergne,
``The Weil Representation, Maslov Index and Theta Series,''
\textit{Birkh\"auser, Boston} (1980).

\bibitem[Mermin1990]{Mermin1990}
N.~D.~Mermin,
``Simple unified form for the major no-hidden-variables theorems,''
\textit{Am. J. Phys.} \textbf{58}, 731--734 (1990).

\bibitem[PaperX]{PaperX}
Z.-L. Chen,
``The Equiangular Characterization of Mermin Pentagrams:
Symplectic Embedding, 10-Ray Cap, and the Failure of the
Collinearity--Obstruction Dictionary,''
\textit{Zenodo, DOI: 10.5281/zenodo.20476659} (2026).

\bibitem[PaperXI]{PaperXI}
Z.-L.~Chen,
``Quadratic Refinement and the Parity of Mermin Pentagrams:\\
The $\beta$-Formula, Geometric Decomposition, and the\\
Kochen--Specker Obstruction from Equiangular Geometry,''
\textit{Zenodo, DOI: 10.5281/zenodo.20482595} (2026).

\bibitem[PaperXIIsuppl]{PaperXIIsuppl}
Z.-L.~Chen,
``Supplementary computational materials for Paper XII,''
\textit{Zenodo, DOI: 10.5281/zenodo.20488394} (2026).

\end{thebibliography}

\end{document}
